\chapter{CLI Reference}
\label{ch:cli-reference}

\section{Invocation}

\begin{lstlisting}[language=bash]
expo-snips <command> [options]
\end{lstlisting}

The binary is declared as \texttt{expo-snips} in
\texttt{packages/cli/package.json} (\texttt{bin}), and the program name and
description are set in \texttt{packages/cli/src/index.ts}:

\begin{lstlisting}[language=JavaScript]
program
  .name("expo-snips")
  .description("Beautiful React Native screens, one command away.")
  .version("0.1.0");
\end{lstlisting}

\textbf{Note.} The version string is hard-coded in the source rather than read from
\texttt{package.json}, so it must be kept in sync manually (see
Chapter~\ref{ch:technical-debt}).

\section{Command summary}

\begin{longtable}{p{2.1cm}p{3.1cm}p{8.6cm}}
\caption{Command summary}\\
\toprule
\textbf{Command} & \textbf{Argument} & \textbf{Purpose} \\
\midrule

\texttt{init} & none & Create a new Expo project using \texttt{create-expo-app}. \\
\texttt{list} & none & List all available screens from the resolved catalog. \\
\texttt{add} & \texttt{[screenId]} & Generate a screen into the current project. \\
\texttt{update} & \texttt{[screenId]} & Update tracked screens to the latest recipe versions. \\
\texttt{status} & none & Show screens tracked by \texttt{.expo-snips.json}. \\
\bottomrule
\end{longtable}

\section{\texttt{expo-snips init}}

\begin{lstlisting}[language=bash]
expo-snips init [--name <name>] [--template <template>]
\end{lstlisting}

\begin{longtable}{p{2.6cm}p{1.6cm}p{10.4cm}}
\caption{\texttt{init} options}\\
\toprule
\textbf{Option} & \textbf{Default} & \textbf{Behaviour} \\
\midrule

\texttt{--name <name>} & unset & Project directory name; when unset the argument
\texttt{.} is used (current directory). \\
\texttt{--template <template>} & \texttt{default} & Passed to
\texttt{create-expo-app} as \texttt{--template <template>}. \\
\bottomrule
\end{longtable}

Behaviour: if \texttt{app.json} already exists in the current working directory the
command prints ``An Expo project already exists in this directory.'' and returns
without doing anything. Otherwise it executes:

\begin{lstlisting}[language=bash]
npx create-expo-app@latest <name|.> --template <template> --yes
\end{lstlisting}

with inherited stdio, then prints next-step guidance. No exit code other than 0 is
set by this command; a failure inside \texttt{create-expo-app} propagates as a
process error from \texttt{execSync}.

\section{\texttt{expo-snips list}}

\begin{lstlisting}[language=bash]
expo-snips list [--registry <url>]
\end{lstlisting}

Resolves the catalog, reports the source (``remote''/``local''), and prints one
line per recipe: the id padded to 24 characters, the display name, and the tags in
square brackets when present.

\begin{lstlisting}[language=bash]
Available screens (24) from local registry:

  app/home                 Home Dashboard Screen [app, home, dashboard]
  auth/login               Login Screen [auth, login, form]
  ...
\end{lstlisting}

Printing an empty catalog prints ``No screens found in the registry.'' and returns
success. See Section~\ref{sec:list-empty} for the discrepancy between this
behaviour and the interactive picker's behaviour.

\section{\texttt{expo-snips add}}

\begin{lstlisting}[language=bash]
expo-snips add [screenId] [options]
\end{lstlisting}

\begin{longtable}{p{3.4cm}p{2.9cm}p{8.5cm}}
\caption{\texttt{add} options}\\
\toprule
\textbf{Option} & \textbf{Value set} & \textbf{Behaviour} \\
\midrule

\texttt{--styling <styling>} & nativewind, stylesheet, tamagui & Explicit styling; rejected when the recipe does not support it. \\
\texttt{--router <router>} & expo-router, react-navigation & Explicit router; rejected when unsupported. \\
\texttt{--ui <ui>} & native, shadcn, custom & Explicit UI primitives; rejected when unsupported. \\
\texttt{--icons <icons>} & lucide, vector-icons, none & Explicit icon library; rejected when unsupported. \\
\texttt{--variant <variant>} & recipe-defined & Screen variant id (for example \texttt{glass}). \\
\texttt{--typescript} & flag & Force TypeScript. \\
\texttt{--javascript} & flag & Force JavaScript (takes precedence over \texttt{--typescript}). \\
\texttt{--no-install} & flag & Skip dependency installation. \\
\texttt{--registry <url>} & URL & Remote registry base; defaults to \texttt{EXPO\_SNIPS\_REGISTRY}. \\
\bottomrule
\end{longtable}

\textbf{Rejected values.} A value that the recipe does not list in
\texttt{supported} aborts with a red error and exit code 1. This applies to
\texttt{--styling}, \texttt{--router}, \texttt{--ui} and \texttt{--icons}.

\textbf{Effective-value notices.} Selecting \texttt{--ui shadcn} prints a note and
overrides \texttt{styling} to \texttt{nativewind} and \texttt{icons} to
\texttt{lucide}, whether or not those were supplied explicitly. A TypeScript
recipe generated with \texttt{--javascript} prints a warning but proceeds.

\subsection*{Examples}

\begin{lstlisting}[language=bash]
# Interactive: pick a screen from a list
expo-snips add

# Named screen, defaults from project detection
expo-snips add auth/login

# Force NativeWind and no dependency installation
expo-snips add app/home --styling nativewind --no-install

# shadcn-style UI (forces nativewind + lucide; adds ui primitives)
expo-snips add auth/register --ui shadcn

# Glassmorphism variant
expo-snips add app/home --variant glass

# JavaScript output
expo-snips add auth/login --javascript

# Point at a locally served registry
expo-snips add app/home --registry http://localhost:3000
\end{lstlisting}

\textbf{Variant availability.} \texttt{--variant glass} has an effect only for
recipes that declare a \texttt{glass} variant: \texttt{app/home},
\texttt{auth/login} and \texttt{onboarding/welcome}. For any other recipe the flag
is accepted but produces default output, because the variant override map is empty.

\section{\texttt{expo-snips update}}

\begin{lstlisting}[language=bash]
expo-snips update [screenId] [--registry <url>] [--yes]
\end{lstlisting}

\begin{longtable}{p{3.4cm}p{2.9cm}p{8.5cm}}
\caption{\texttt{update} options}\\
\toprule
\textbf{Option} & \textbf{Value set} & \textbf{Behaviour} \\
\midrule

\texttt{[screenId]} & recipe id & Update only this recipe; error when no tracked
entry matches. \\
\texttt{--registry <url>} & URL & Remote registry base. \\
\texttt{--yes} & flag & Apply without the confirmation prompt. \\
\bottomrule
\end{longtable}

Behaviour:

\begin{itemize}
  \item No tracked screens: prints guidance to use \texttt{add} and returns.
  \item Nothing changed: prints ``All tracked screens are up to date.''
  \item Changes found: prints the file count, a coloured diff per changed file, then
        a confirmation prompt defaulting to \emph{no}; declining prints ``Update
        cancelled.''
  \item Applied: rewrites the files and prints the number of updated screens.
  \item A tracked recipe missing from the registry is skipped with a warning.
\end{itemize}

\begin{lstlisting}[language=bash]
expo-snips update
expo-snips update auth/login
expo-snips update --yes
\end{lstlisting}

\textbf{Important behaviour.} \texttt{update} compares the newly rendered content
with the file currently on disk. If the developer has edited a generated file, the
update will overwrite those edits once applied. The diff is the only safeguard.

\section{\texttt{expo-snips status}}

\begin{lstlisting}[language=bash]
expo-snips status
\end{lstlisting}

Reads \texttt{.expo-snips.json} from the current working directory and prints each
entry as \texttt{<recipeId padded to 24> v<recipeVersion>} with an optional
\texttt{(variant: <id>)} suffix, followed by a hint to run
\texttt{expo-snips update}.

\begin{lstlisting}[language=bash]
Tracked screens (2):

  auth/login               v0.3.0
  app/home                 v0.1.0 (variant: glass)

Run expo-snips update to update them to the latest version.
\end{lstlisting}

With no tracked screens it prints ``No screens tracked in this project yet.''

\section{Exit codes and error messages}

\begin{longtable}{p{4.4cm}p{9.6cm}}
\caption{Failure conditions and messages}\\
\toprule
\textbf{Condition} & \textbf{Message and exit code} \\
\midrule

Interactive selection cancelled & no message; exit code 1. \\
Screen id not found (\texttt{add}) & ``Screen not found: <id>'' plus a pointer to
\texttt{list}; exit code 1. \\
Option unsupported by recipe & ``Recipe "..." does not support <option> "...".''
exit code 1. \\
No files matched the options & ``No files matched the selected options. Nothing
generated.''; exit code 1. \\
No tracked entries for a given \texttt{update <id>} & ``No installed screens match
"<id>".'' plus pointer to \texttt{status}; exit code 1. \\
Empty catalog during interactive \texttt{add} & guidance to reinstall the CLI or
configure a registry; exit code 1. \\
Empty catalog during \texttt{list} & ``No screens found in the registry.''; exit code 0. \\
\bottomrule
\end{longtable}

\label{sec:list-empty}

\section{Version discrepancy note}

The CLI's \texttt{--version} output is the literal string \texttt{0.1.0} defined in
\texttt{src/index.ts}. The \texttt{package.json} also declares \texttt{0.1.0}, so
they currently agree; however, changing one without the other would silently
diverge. The VS Code extension avoids this class of problem by letting the editor
read the version from its manifest.
